% !TeX root = ../papers/nc-p3.tex
\section{整二次錐の全商と平面成分からの正規一次直線}
\label{sec:quadric-cone-quotients}
$D$ を階数3の二次式による四変数多項式環の商とする。
\begin{theorem}\label{thm:quadric-cone-quotients}
$\check D$ への一次で同型な全射を持つ標準14641型には、次の明示的な全同型類がある。
$\epsilon^2=0$ とし、三つの欄は二次双対への追加関係 $a_i$ を表す。
\begin{enumerate}
\item $a_i=I+\epsilon N$。$N$ は零、非零冪零Jordan行列、$\operatorname{diag}(1,-1)$ の三種。
\item $a_i=P$ は非スカラー可逆二次行列。共通倍率と共役を除き、
異なる根の比 $u\in k^\times\setminus\{1\}$ ($u\sim u^{-1}$) または非自明Jordan行列を取る。
\item $P=\operatorname{diag}(1,-1)$, $a_i=(1+(-1)^i\epsilon)P$ の一種。
\end{enumerate}
全て最小周期1であり、これら三欄の重複はない。第一欄の $N=0$ は可換多項式窓である。
最後の一種だけが余分な部分を持たず二次錐そのものを点対応とする。
整二次錐商を持つ有限標準Frobenius候補では、このリストに入ることもKoszul性も自動である。
\end{theorem}
\begin{proof}[退化Clifford環]
$T=k[\epsilon]/(\epsilon^2)$ とし $M_2(T)$ に $j^2=1$, $jBj=\bar B$
（$\epsilon$ の符号を変える）で奇元 $j$ を加える。
一次文字の像を $(U+\epsilon uI)j$, $U\in\mathfrak{sl}_2(k),u\in k$ とする。
平方は $U^2=-\det(U)I$ なので、これは階数3の形式のClifford表示である。
中心二次式 $h$ を独立にした前節の $k[h]$ 自由基底は退化した形式でも有効である。
例えば対角化して平方を $c_i h$、異なる二文字を反交換させれば順序語が基底となる。
$h=1$ では二文字積は $UV+\epsilon(uV-vU)$ となり、その線型包は
$H=\{P+\epsilon Q:\operatorname{tr} Q=0\}$ の全七次元である。
三文字は $M_2(T)j$、四文字は $M_2(T)$ の全空間を張る。これは無跡行列の積が
$M_2(k)$ を張ることと可逆な無跡文字を用いて分かる。
従って各斉次片のこの表示は単射である（次数4以後は $h$ による移送を使う）。

$M_2(T)$ のFrobenius汎関数を $\operatorname{tr}$ の $\epsilon$ 係数に取ると $H^\perp=kI$ となる。
四次traceに対応する行列を $b_i$ とする。前節と同じ二次対合から
$\{a\in H:b_i a\in kI\}=ka_i$ である。$b_i$ が非可逆ならその全行列空間での左乗法核の
$k$ 次元は2以下である。$T$ 上の行列は単元行列の左右乗法でSmith形にでき、
核の次元は $T^2$ 上の核の次元の二倍である。従って非可逆な場合のSmith形は
$\operatorname{diag}(1,\epsilon)$ だけである。その核にある $a_i$ は
$\epsilon N$ で $N\ne0$ は跡零かつ階数1、特に各 $b_i$ の剰余は非零である。
しかし $C_1a_i C_1$ は $\epsilon M_2(k)$ の全体を張る。実際、$N$ を行列単位 $E_{12}$ に
共役すれば $U E_{12} V$ で無跡な $U,V$ の第一列と第二行を任意に取れる。
前位置のtraceがこの空間を消すなら $b_{i-1}$ の剰余は零となり矛盾する。
従って全 $b_i,a_i$ が可逆である。
\end{proof}

\begin{proof}[移送、同型、正則性]
三次核次元4より再び $a_i C_1=C_1 a_{i+1}$ である。
$a_i=P_i+\epsilon Q_i$ と書く。全無跡 $U$ で $P_i U P_{i+1}^{-1}$ が無跡なので
$P_i$ は射影的に同じであり、$P_i=P$ とする。$D_i=Q_iP^{-1}$ と置くと
$a_i(U+\epsilon uI)\bar a_{i+1}^{-1}$ の $\epsilon$ 係数の条件は、
全無跡 $S=P U P^{-1}$ で $D_i S+S D_{i+1}$ がスカラーであることに等しい。
$D_i=d_i I+N_i$ と跡と無跡に分ける。無跡二行列の反交換子はスカラーなので、
無跡部分で $(d_i+d_{i+1})S+[N_i-N_{i+1},S]/2=0$ である。
無跡空間上の写像の跡から $d_{i+1}=-d_i$、中心化条件から $N_{i+1}=N_i$ となる。
よって $D_i=N+(-1)^i dI$, $\operatorname{tr}(NP)=0$, $d\operatorname{tr}P=0$ で尽くされる。

二次錐の射影自己同型は、全体倍率を除いて $(U,u)\mapsto(GUG^{-1},c u+f(U))$、
$c\ne0$, $G\in\mathrm{PGL}_2$ である。これは三次元の非退化二次形式の射影直交群と
一次元の根基を使えば従う。切断グラフの同型が抽象的に与えられても超平面束は保たれる。
実際、二次錐の極小解消 $\mathbb F_2$ で例外切断を $E$、母線を $F$ とすると、
引戻し超平面類は $E+2F$ であり、自己同型は唯一の負切断と母線を保つ。
$u$ の平行移動 $f(U)=\operatorname{tr}(K U)$ ($\operatorname{tr}K=0$) は
行列環の $I+\epsilon K$ による共役で実現する。奇文字では
$\{K,U\}=\operatorname{tr}(KU)I$、偶文字では $Q_i\mapsto Q_i+[K,P]$ となる。
従って $N$ は $(1-\operatorname{Ad}P)\mathfrak{sl}_2$ の加算まででよい。
非スカラー $P$ ではこの像は二次元で、$\operatorname{tr}(NP)=0$ の平面に等しい。
$P$ がスカラーなら $d=0$、残る $N$ は共役と非零倍率で割る。
これでリストの全同型を得る。交互な $d$ の符号は $c=-1$ で逆転できる。

逆行列の $\epsilon$ 係数にも跡がない。実際 $\operatorname{tr}(D_iP^{-1})=
(\operatorname{tr}P\operatorname{tr}D_i-\operatorname{tr}(D_iP))/\det P=0$。
従って各逆行列は二次の補元に持ち上がり、元との積は $h^2$ となる。
前節の正規二次式商のKoszul性と $C/(h^2)$ でのFrobenius対合を同様に使える。
$h$ を消すと四文字の外積環である点も変わらない。

可換表示以外では、前節と同じ $p\wedge y$ の議論で唯一の二次元既約成分が
二次錐である。交代部分が階数2でその台の直線が曲面に含まれなければ、
その直線の一般点に別の点対応が出る。含まれれば錐の滑らかな点で、前節の接計算から
グラフより大きな接空間が出る。交代部分が階数4のときだけグラフがそのまま得られる。
Clifford表示では交代二形式を、三つの無跡文字間の交換子と、根基文字とそれらとの交換子に分ける。
前者は $P$ の無跡部分 $P_0$、後者は $Q_i$ とそれぞれ無跡行列空間で同一視される。
両者を外積して非零となる条件は $\operatorname{tr}(P_0 Q_i)\ne0$ である。
実際、三無跡文字の交換子は外積二乗と無跡空間の同型で、
その体積対合は $\operatorname{tr}([U,V]W)$ の非零倍である。
リストの最後に限りこの跡が非零である。
\end{proof}

\begin{theorem}\label{thm:plane-surface-normal}
標準窓の全辺の射影ファイバーが空か一点で、一辺の射影像に平面成分があるとする。
この平面から正則正規一次直線列が復元され、中央化して得る三生成商は
三変数多項式環の次数付きtwistである。従って零次元射影ファイバーで二次元成分を持つ全標準窓は、
滑らかな二次曲面商、整二次錐商、または正規一次直線列の平面商のいずれかに入る。
\end{theorem}
\begin{proof}
前節の共有像の同型と次数上界から、平面の次の像も平面である。平面と同型な
埋込み曲面は次数2の整二次曲面になれない。射影像が平面三空間全体なら
自己同型は線型で、六関係はそのグラフの全関係となり、可換多項式窓に帰着する。
以下、平面の列に沿って座標を運び、商を三変数の可換多項式環 $S$ の反復に取る。
一次核を $kt_i$ とする。

元の標準の単純加群の右線型解消に左から作用されるこの商をテンソルする。
始点を固定して得る複体は $S$ 上の長さ4の有限自由複体で、
平面上の各点では順に各微分の階数が $1,3,3,1$ である。
中央二写像は隣接関係を片方の点で縮約したものと相補関係の転置であり、
グラフ上で階数3、両端は点の非零評価である。従って射影平面上では完全で、
各ホモロジーは無関係イデアルに台を持つ有限長 $S$ 加群である。
$\operatorname{Ext}^p_S(-,S)$ は非零な有限長加群上 $p=3$ にだけ残る。
複体のhyper-Extの全次数 $p+q$ は4以下なので $q>1$ のホモロジーは零である。
第零は $k$、Euler級数は $(1-z)^4/(1-z)^3=1-z$ だから第一は $k(-1)$ である。
この計算は各単純に対する $\operatorname{Tor}^A(k_i,S)$ を内部次数ごとに与える。
各終点を固定した $S$ の左極小解消を見ると、核はただ一つの一次元 $t$ の右乗法による
自由加群であり、その乗法は単射である。反対側でも同じ計算を行うと、
商への核は両方からその一次元だけで生成され、左右に正則となる。
定理\ref{thm:regular-normal-linear}による一移動で中央化すると、商の点対応は
依然平面のグラフである。その移動は射影線型なので、三生成三関係の商は多項式環の次数付きtwistである。
最後の被覆は真の曲面因子の総次数が高々2であることから従う。
\end{proof}
